% ---------------------------------------------------------------------------------------
% MODELO DE TCC - INSTITUTO FEDERAL DE SERGIPE (IFS)
% ---------------------------------------------------------------------------------------
% Autor do Modelo: Mikael dos Santos Gois (mikaelgois1@gmail.com)
% Baseado nos modelos de: 
%   - Diogo D. Moreira (TCC, Dissertação e Tese)
%   - DECOM/UFOP (Monografia de Ciência da Computação)
% Modelo Base: abnTeX2 adaptado para as normativas internas de TCC do IFS (Art. 31 e 34)
% Versão: 1.0 (2026)
% ---------------------------------------------------------------------------------------
% Este arquivo faz parte de um modelo opensource para estudantes do IFS.
% Não remova os créditos dos autores originais e do modelo base.
% Este modelo é fornecido "no estado em que se encontra", sem garantias de qualquer tipo.
% ---------------------------------------------------------------------------------------

\begin{resumo}
  A programação é uma competência fundamental no século XXI, mas o ensino introdutório enfrenta barreiras
  significativas que desmotivam o aprendizado inicial, como a barreira da língua inglesa e a complexidade
  sintática prematura. Este trabalho apresenta o projeto Luso, uma linguagem de programação educacional
  projetada com sintaxe em português. Baseado nas teorias de Localização de Software e Programação Gradual,
  o Luso visa reduzir a carga cognitiva de estudantes novatos, servindo como um andaime (scaffolding)
  para a transição futura para linguagens profissionais como Java e Python. O projeto propõe um ecossistema
  que inclui um transpilador, uma IDE e suporte a ferramentas de desenvolvimento terceiras, buscando
  alinhar a tecnologia à didática do professor. Além disso, o projeto conta com documentação educacional
  para usuários e documentação técnica para contribuidores do projeto.

  \vspace{\onelineskip}

  \noindent
  \textbf{Palavras-chave}: Ensino de Programação; Linguagens Graduais; Localização de Software; Barreira
  de Sintaxe; Carga Cognitiva.
\end{resumo}